En mi caso muestro el panel de monitoreo en tiempo real de SGIRI, pensado para que un operador supervise el estado de la red desde una sola pantalla.
Como se ve en la captura, el panel concentra cuatro métricas principales: latencia promedio, paquetes UDP con su porcentaje de pérdida, conexiones TCP activas y el estado general del servicio.
Debajo se actualiza cada segundo una gráfica de latencia junto con una tabla de eventos recientes que identifica el nodo, el protocolo y si hubo pérdida.
Aquí el WebSocket cumple el papel central: el servidor empuja las mediciones y alertas apenas ocurren, y el navegador las dibuja al vuelo sin tener que hacer peticiones HTTP repetitivas.
Si la conexión se cae, el indicador de estado avisa al operador, los eventos registran la pérdida y el cliente intenta reconectarse automáticamente para retomar la lectura apenas vuelve el enlace.